\section{Introduction}

Deep forecasters trained with mean squared error produce predictions that are visibly smoother than the series they imitate: a well-tuned PatchTST on ETTh1 predicts with an amplitude ratio of 0.84 relative to the realized series, and half of its point predictions fall strictly between a level reference and the realized value. This signature is widely treated as a defect: frequency-domain losses attribute it to label autocorrelation \citep{fredf}, auxiliary classification heads to the failure of squared losses on high-entropy features \citep{hcan}, and shape-aware objectives penalize it directly \citep{dilate}---all premised on the smoothness being pathological and repair improving accuracy.

This premise needs qualification. A squared-loss-optimal forecast is the conditional mean of the future given the context, which is necessarily smoother than any realization, because realizations contain unpredictable noise that no point forecast should imitate: the over-smoothing signature is a property of the optimal predictor itself, not only of imperfect models. Whether a trained model is \emph{too} smooth depends on how its signature compares with the signature an oracle would leave on the same data. This reframing decomposes over-smoothing into \emph{rational shrinkage}, which an optimal predictor must keep, and \emph{unexploited structure}, which can in principle be repaired.

We make the decomposition measurable. Our diagnostic, the \emph{oracle gap}, bootstraps a trained model's residuals to synthesize alternative realizations of the future and compares the amplitude signature the model's predictions would exhibit as an oracle against them with the signature observed against the true data. The gap needs no training and runs in seconds per cell: positive indicates residuals with structure the model failed to extract, non-positive indicates rational shrinkage.

The diagnostic is predictive, not just descriptive: across fifteen dataset-backbone cells, the sign of the gap predicts the sign of the intervention's effect in all but one cell (14/15, $p < 5 \times 10^{-4}$ under a sign-symmetric null); the exception is diagnosable in advance from the anchor geometry. A separate 665-cell audit of thirteen foundation models and a trained PatchTST on fifty GIFT-Eval variants \citep{gifteval} corroborates from the accuracy side, and settles several folk questions: ETT is saturated, transit and energy series retain exploitable structure across all thirteen models, scale is irrelevant from 8M to 311M parameters, and accuracy rank anti-correlates with leftover structure.

Guided by the diagnostic, we introduce \rfc, a region-weighted loss that increases the penalty where the prediction under-shoots the realized deviation from a level reference while leaving the zero-loss set of MSE untouched; a Gaussian analysis bounds the induced fixed-point shift at roughly $0.4\alpha\sigma$, so the loss acts as a calibrated counter-bias rather than a new objective. On Electricity, \rfc~reduces MSE by 2.3\% across three seeds and horizons from 96 to 720 steps; on non-positive-gap datasets the diagnostic withholds the repair, avoiding the harm it would otherwise cause.

Where shrinkage is rational, a model can still be made more useful by knowing when it is likely to fail. A light-weight descriptor head predicts structural properties of the future window, with labels computed for free from training data. Such heads cannot repair point accuracy at the scales we study, echoing scale-dependent findings for multi-token prediction \citep{gloeckle2024mtp,fsf}, but the predicted change-point probability is a reliable forward-looking risk signal: abstaining on the 20\% highest-risk windows reduces the remaining MSE by up to 5.3\% on task models and 9.6\% on a fine-tuned foundation model, outperforming oracle change-point rankings and static baselines. Predicting future structure is more useful for deciding when to trust a forecast than for improving it.

\textbf{Contributions.}
\begin{itemize}\setlength\itemsep{1pt}
    \item A two-component theory of over-smoothing---rational shrinkage versus unexploited structure---with a conditional improvement interval $(0, 2b)$ and a Jensen-gap argument for supervising structure through auxiliary heads.
    \item The oracle gap, a training-free diagnostic of repairable structure: its sign predicts intervention outcomes in 14 of 15 cells (the exception has an identified geometric cause), corroborated by a 665-cell audit.
    \item \rfc, a region-weighted loss that repairs over-smoothing where the diagnostic is positive and is withheld where it would harm, with a bounded fixed-point shift.
    \item An audit of forecasting benchmarks through repairable structure: saturation evidence for ETT, positive cells in transit and energy series, and flat scaling across model size.
    \item Evidence that predicted structural descriptors are more valuable at decision time than at training time, including a selective-prediction gain of 5.3\% exceeding oracle change-point ranking.
\end{itemize}
